\section*{Capítulo \ref{ch:g12:ka-and-pka} --- Ácidos e bases: Ka e pKa}

\begin{solution}{exo:g12:ka-and-pka:1}
Bases conjugadas: \ce{HCOO-}, \ce{NH3}, \ce{OH-}, \ce{CO3^{2-}}. Ácidos
conjugados: \ce{NH4+}, \ce{H2O}, \ce{H3O+}, \ce{CO2,H2O}. \omterm{def:g12:ka-and-pka:oxonium}{Espécies
anfóteras}: \ce{H2O} e \ce{HCO3-}.
\end{solution}

\begin{solution}{exo:g12:ka-and-pka:2}
\omterm{def:g9:acids-bases-ph:ph}{pH} 1.0; 3.0; $-\log(\num{5.0e-3}) = 2.30$.
\end{solution}

\begin{solution}{exo:g12:ka-and-pka:3}
$10^{-4.5} = \qty{3.2e-5}{mol/L}$. A \omterm{def:g9:acids-bases-ph:ph}{pH} 11.2:
$[\ce{H3O+}] = \qty{6.3e-12}{mol/L}$ e
$[\ce{OH-}] = \num{1.0e-14}/\num{6.3e-12} = \qty{1.6e-3}{mol/L}$.
\end{solution}

\begin{solution}{exo:g12:ka-and-pka:4}
\ce{HCOOH + H2O <=> HCOO- + H3O+},
$K_a = \dfrac{[\ce{HCOO-}][\ce{H3O+}]}{[\ce{HCOOH}]\,c^\circ}$;
\ce{NH4+ + H2O <=> NH3 + H3O+},
$K_a = \dfrac{[\ce{NH3}][\ce{H3O+}]}{[\ce{NH4+}]\,c^\circ}$.
\end{solution}

\begin{solution}{exo:g12:ka-and-pka:5}
O ácido metanoico (menor $pK_a$). A base mais forte é o \ce{CH3COO-}, a
base do ácido mais fraco.
\end{solution}

\begin{solution}{exo:g12:ka-and-pka:6}
$[\ce{A-}] = [\ce{H3O+}] = \qty{1.0e-3}{mol/L}$;
$[\ce{HA}] = 0.050 - 0.001 = \qty{0.049}{mol/L}$;
$K_a = \num{1.0e-6}/0.049 = \num{2.0e-5}$; $pK_a = 4.69$.
\end{solution}

\begin{solution}{exo:g12:ka-and-pka:7}
$K_a = \num{6.31e-5}$; $x^2 + \num{6.31e-5}\,x - \num{1.26e-6} = 0$ dá
$x = \qty{1.09e-3}{mol/L}$: \omterm{def:g9:acids-bases-ph:ph}{pH} 2.96.
\end{solution}

\begin{solution}{exo:g12:ka-and-pka:8}
$\mathrm{pH} = 14.0 + \log(\num{2.0e-3}) = 14.0 - 2.70 = 11.3$.
\end{solution}

\begin{solution}{exo:g12:ka-and-pka:9}
\ce{HF}, \ce{HCOOH}, ácido ascórbico, \ce{C6H5COOH}, \ce{CH3COOH},
\ce{CO2,H2O}, \ce{NH4+}, \ce{HCO3-}. Um \omterm{def:g12:ka-and-pka:strong-acid}{ácido forte} fica bem abaixo,
além do fim da escala: ele reage totalmente com a água.
\end{solution}

\begin{solution}{exo:g12:ka-and-pka:10}
$[\ce{H3O+}] = 10^{-5.6} = \qty{2.5e-6}{mol/L}$, 25 vezes mais que na
água pura (\qty{1.0e-7}{mol/L}).
\end{solution}

\begin{solution}{exo:g12:ka-and-pka:11}
$K_a K_b = \dfrac{[\ce{NH3}][\ce{H3O+}]}{[\ce{NH4+}]c^\circ} \times
\dfrac{[\ce{NH4+}][\ce{OH-}]}{[\ce{NH3}]c^\circ} = K_e$, logo
$pK_a = pK_e - pK_b = 14.0 - 4.74 = 9.26$.
\end{solution}

\begin{solution}{exo:g12:ka-and-pka:12}
A \qty{0.0010}{mol/L}: $x^2 + \num{1.74e-5}\,x - \num{1.74e-8} = 0$,
$x = \qty{1.24e-4}{mol/L}$, \omterm{def:g9:acids-bases-ph:ph}{pH} 3.91: uma subida de 0.52, não de 1.
Diluído, um \omterm{def:g12:ka-and-pka:strong-acid}{ácido fraco} reage numa proporção maior com a água (aqui
\qty{12}{\percent} contra \qty{4}{\percent}), o que compensa em parte a
\omterm{def:g10:concentration-and-dilution:dilution}{diluição}.
\end{solution}

\begin{solution}{exo:g12:ka-and-pka:13}
Como ácido: \ce{HCO3- + H2O <=> CO3^{2-} + H3O+}, $K = 10^{-10.33}$. Como
base: \ce{HCO3- + H2O <=> H2CO3 + OH-} (com \ce{H2CO3} representando \ce{CO2,H2O}), $K = K_e/K_a = 10^{-14.0+6.37}
= 10^{-7.63}$. A reação como base tem a maior constante: a solução é
levemente básica.
\end{solution}

\begin{solution}{exo:g12:ka-and-pka:14}
Um ácido, por mais diluído que seja, não pode tornar uma \omterm{def:g9:acids-bases-ph:acidic}{solução básica}.
Quando os \omterm{def:g12:ka-and-pka:oxonium}{íons oxônio} do próprio ácido ficam menos numerosos que os
\qty{1e-7}{mol/L} fornecidos pela água, os \omterm{def:g9:ions:ion}{íons} da água predominam: o \omterm{def:g9:acids-bases-ph:ph}{pH}
se aproxima de 7 por baixo e nunca passa disso.
\end{solution}

\begin{solution}{exo:g12:ka-and-pka:15}
$n = 0.500 / 176.0 = \qty{2.84e-3}{mol}$, $c = \qty{0.0142}{mol/L}$;
$K_a = \num{9.12e-5}$; $x^2 + \num{9.12e-5}\,x - \num{1.30e-6} = 0$ dá
$x = \qty{1.09e-3}{mol/L}$: \omterm{def:g9:acids-bases-ph:ph}{pH} 2.96.
\end{solution}

\begin{solution}{pb:g12:ka-and-pka:1}
\textbf{1.} \ce{H-C(=O)-O-H}: o hidrogênio cedido é o do \ce{O-H}.

\textbf{2.} \ce{HCOOH}/\ce{HCOO-}; \ce{HCOOH + H2O <=> HCOO- + H3O+}.

\textbf{3.} $K_a = \dfrac{[\ce{HCOO-}][\ce{H3O+}]}{[\ce{HCOOH}]c^\circ} =
10^{-3.74} = \num{1.82e-4}$.

\textbf{4.} Abaixo do ácido etanoico na escala ($3.74 < 4.76$): mais forte.

\textbf{5.} $M = \qty{46.0}{g/mol}$: $0.010 \times 46.0 = \qty{0.46}{g}$.

\textbf{6.} \ce{HCOOH}: $0.010 - x$; \ce{HCOO-}: $x$; \ce{H3O+}: $x$.

\textbf{7.} $\dfrac{x^2}{0.010 - x} = \num{1.82e-4}$.

\textbf{8.} $x^2 + \num{1.82e-4}\,x - \num{1.82e-6} = 0$:
$x = \qty{1.26e-3}{mol/L}$.

\textbf{9.} $\mathrm{pH} = -\log(\num{1.26e-3}) = 2.90$.

\textbf{10.} $\num{1.26e-3}/0.010 = \qty{12.6}{\percent}$.

\textbf{11.} \qty{1.26e-3}{mol/L} é mais de dez mil vezes
\qty{1e-7}{mol/L}: os \omterm{def:g9:ions:ion}{íons} da própria água são desprezíveis.

\textbf{12.} \omterm{def:g9:acids-bases-ph:ph}{pH} 2.0.

\textbf{13.} $0.010 / \num{1.26e-3} \approx 8$ vezes mais.

\textbf{14.} O ácido clorídrico é um \omterm{def:g12:ka-and-pka:strong-acid}{ácido forte}, inteiramente
transformado em \omterm{def:g9:ions:ion}{íons}; o ácido metanoico é fraco: a sua reação com a água
para num equilíbrio em que a maior parte dele continua como \ce{HCOOH}.

\textbf{15.} \ce{HCOOH + HCO3- -> HCOO- + CO2 + H2O}.

\textbf{16.} Multiplicando o numerador e o denominador de
\[
  K = \frac{[\ce{HCOO-}]\,[\ce{CO2}]}{[\ce{HCOOH}]\,[\ce{HCO3-}]}
\]
por $[\ce{H3O+}]$, obtém-se $K = K_a(\ce{HCOOH}) / K_a(\ce{CO2,H2O}) =
\num{1.82e-4}/\num{4.3e-7} \approx 420$.

\textbf{17.} Sim, $K \gg 1$; o dióxido de carbono borbulha.

\textbf{18.} O hidróxido é uma \omterm{def:g12:ka-and-pka:strong-acid}{base forte} e atacaria ele mesmo a pele; o
hidrogenocarbonato é uma \omterm{def:g12:ka-and-pka:strong-acid}{base fraca} que neutraliza o ácido e é suave
mesmo em excesso.

\textbf{19.} \omterm{def:g9:acids-bases-ph:ph}{pH} 2.9.
\end{solution}
